\begin{abstract}
Este trabajo presenta una implementación personalizada de U-Net en PyTorch para segmentación binaria de vasos retinianos, desarrollada sin usar bibliotecas de segmentación preconstruidas. Se evaluaron decisiones de entrenamiento y arquitectura mediante un estudio de ablación, comparando funciones de pérdida (BCE, Dice y combinada), capacidad del modelo y variantes con conexiones de salto enriquecidas. El modelo se entrenó en DRIVE y se evaluó tanto en distribución (DRIVE) como fuera de distribución (CHASE\_DB1), reportando sensibilidad, especificidad, F1 y AUC-ROC. Además, se analizó la brecha de generalización entre dominios y se estudió una estrategia de adaptación basada en preprocesamiento (CLAHE + normalización). Finalmente, se realizó un análisis cualitativo de fallos para distinguir el comportamiento en capilares finos frente a vasos grandes y explicar sus causas arquitectónicas.
\end{abstract}
